\chapter{Delimitación del Objeto de Estudio}
\label{chap:delimitacion}

\noindent Este capítulo delimita el objeto de estudio de la investigación mediante la identificación del problema, la formulación de la pregunta de investigación, el planteamiento de hipótesis y el establecimiento de objetivos específicos.

% ============================================================================
\section{Problema de Investigación}
\label{sec:problema_investigacion}
% ============================================================================

La evaluación precisa del comportamiento sedentario constituye un pilar fundamental para la estratificación del riesgo y el diseño de intervenciones efectivas en salud pública. Evidencia reciente de meta-análisis a gran escala confirma su asociación independiente con un mayor riesgo de mortalidad por todas las causas, incidencia de enfermedades cardiovasculares y diabetes mellitus tipo 2 \cite{Katzmarzyk2023Impact,Santos2023CostInaction}. Sin embargo, los métodos de evaluación convencionales presentan una dicotomía metodológica que limita su aplicabilidad y precisión. Los instrumentos subjetivos, como los cuestionarios de autoinforme, aunque escalables, adolecen de sesgos de memoria y deseabilidad social que comprometen su validez \cite{Healy2024,Prince2008}. En contraparte, las técnicas objetivas de referencia, como la calorimetría indirecta o la actigrafía de grado investigativo, se restringen a entornos de laboratorio que no capturan la complejidad del comportamiento en condiciones de vida libre \cite{Migueles2022GRANADA}.

La ubicuidad de los dispositivos portátiles de consumo ha generado una oportunidad sin precedentes para la monitorización objetiva, continua y longitudinal del comportamiento en el entorno ecológico del individuo. No obstante, la transición de la recolección masiva de datos a la generación de conocimiento clínicamente accionable presenta un desafío metodológico significativo. El reto principal reside en traducir flujos de datos multivariados y heterogéneos ---actividad física, patrones de frecuencia cardíaca, variabilidad de la frecuencia cardíaca (\textit{HRV}) y gasto energético--- en una clasificación del sedentarismo que sea fisiológicamente significativa e individualizada \cite{Deka2023Nonlinear}. Los modelos de inteligencia artificial basados en aprendizaje profundo, comúnmente denominados ``de caja negra'', aunque han demostrado una alta precisión en la clasificación de actividades, a menudo carecen de la interpretabilidad necesaria para su adopción en la práctica clínica, donde la auditabilidad de una recomendación es un imperativo ético y profesional \cite{Vellido2020Importance}.

El análisis de la literatura revela brechas metodológicas significativas que justifican la presente investigación. A pesar de la existencia de un precedente metodológico que combina \textit{clustering} \textit{K-Means} con sistemas de inferencia difusa \textit{Mamdani}, no se identifica ningún estudio que haya validado este enfoque híbrido mediante protocolos de \textit{Leave-One-User-Out} (\textit{LOUO}), estrategia esencial para evaluar la generalización inter-usuario en aplicaciones clínicas \cite{Rehman2024LOSO,Mathew2024LOSO,Bassani2025DLHAR}. Adicionalmente, la mayoría de los estudios de clasificación de sedentarismo se basan exclusivamente en señales de movimiento (acelerómetro, giroscopio), abordando el problema desde una perspectiva postural en lugar de fisiológica, y no se identifica ningún estudio que haya utilizado \textit{HRV} como variable de entrada para clasificar comportamiento sedentario, a pesar de su potencial para diferenciar estados fisiológicos distintos \cite{Godkin2025Context,Marino2024ARIC}. Asimismo, no existe evidencia de sistemas de lógica difusa aplicados a datos de \textit{Apple Watch} (acelerómetro + fotopletismografía) para la clasificación de sedentarismo, ni estudios que hayan utilizado datos recolectados durante múltiples años en condiciones de vida libre sin restricciones, limitando la validez ecológica de los modelos desarrollados \cite{Fuller2021Predicting}. Finalmente, no se identifica ningún estudio que haya desarrollado y validado un sistema de clasificación con $N=10$ o menos, un escenario común en estudios piloto y fases iniciales de desarrollo \cite{Bolger2013}.

Este vacío metodológico ---la ausencia de un sistema que integre de forma sinérgica múltiples biomarcadores de vida libre en una clasificación del sedentarismo que sea robusta, empíricamente validada e interpretable--- constituye el núcleo del problema que se aborda en esta investigación. Los análisis estadísticos univariados han resultado insuficientes para capturar las complejas interacciones no lineales entre la actividad física, la función autonómica y la respuesta metabólica \cite{Aubert2022GlobalMatrix}. Por ejemplo, la relación entre el \textit{HRV} y la actividad física es intrínsecamente no lineal y dependiente del contexto, lo que requiere enfoques de modelado que puedan manejar la incertidumbre y la vaguedad inherentes a los sistemas biológicos \cite{Deka2023Nonlinear}.

Para resolver este problema, esta investigación desarrolló un sistema de evaluación basado en un sistema de inferencia difusa tipo \textit{Mamdani}. Este enfoque se apartó de una hipótesis inicial centrada en correlacionar datos objetivos con percepciones subjetivas, para en su lugar, validar el sistema experto contra una ``verdad operativa'' derivada de los propios datos mediante un análisis de conglomerados no supervisado. De este modo, el estudio se enfocó en resolver el problema fundamental de crear un clasificador objetivo, interpretable y validado para el comportamiento sedentario en condiciones de vida libre. El clasificador utiliza datos longitudinales recopilados a través de dispositivos de consumo.

% ============================================================================
\section{Pregunta de Investigación}
\label{sec:pregunta_investigacion}
% ============================================================================

¿De qué manera un sistema basado en lógica difusa puede clasificar el comportamiento sedentario a partir de datos biométricos recopilados en condiciones de vida libre por dispositivos ``\textit{wearables}''?

% ============================================================================
\section{Hipótesis}
\label{sec:hipotesis}
% ============================================================================

\subsection{Hipótesis Conceptual (HC)}
\label{subsec:hipotesis_conceptual}

La clasificación del comportamiento sedentario generada por el sistema de inferencia difusa, basado en reglas que integran biomarcadores de actividad y función cardiovascular, exhibe una alta concordancia con la clasificación objetiva obtenida mediante un análisis de conglomerados no supervisado sobre los mismos datos.

\subsection{Hipótesis Nula ($H_0$)}
\label{subsec:hipotesis_nula}

No existe una concordancia estadísticamente significativa entre la clasificación del comportamiento sedentario generada por el sistema de inferencia difusa y la clasificación obtenida mediante el análisis de conglomerados. Cualquier concordancia observada no es superior a la que se obtendría por azar.

% ============================================================================
\section{Objetivo General}
\label{sec:objetivo_general}
% ============================================================================

Construir y validar un sistema interpretable, basado en lógica difusa, para la clasificación del comportamiento sedentario a partir de datos biométricos de \textit{wearables} recopilados en condiciones de vida libre.

% ============================================================================
\section{Objetivos Específicos}
\label{sec:objetivos_especificos}
% ============================================================================

\begin{enumerate}
    \item Analizar los datos biométricos longitudinales para derivar un conjunto de variables semanales que representen cuantitativamente los patrones de actividad física y función cardiovascular de la cohorte de estudio.
    
    \item Identificar los perfiles de comportamiento sedentario inherentes a los datos mediante la aplicación de un algoritmo de agrupamiento no supervisado, para establecer una clasificación de referencia empírica.
    
    \item Diseñar un sistema de inferencia difusa que modele la relación entre las variables biométricas y los perfiles de comportamiento, el cual establezca reglas lingüísticas con fundamento fisiológico.
    
    \item Evaluar el desempeño del sistema de inferencia difusa mediante la determinación de su nivel de concordancia con los perfiles de comportamiento identificados por el método de agrupamiento.
    
    \item Examinar la contribución de los componentes del sistema a través de un análisis de sensibilidad para determinar la importancia relativa de sus variables de entrada en el rendimiento clasificatorio.
\end{enumerate}
